%% ch03-configuration-ref.tex — Configuration Reference
%% JA4proxy Technical Reference Manual

\chapter{Configuration Reference}
\index{configuration}
\label{ch:configuration}

\begin{chapterquote}
Every behavioural parameter in \ja\ is configurable in \code{config/proxy.yml}.
All keys except listen port, Redis URL, and TLS certificate paths support
hot-reload via \code{SIGHUP} or the Management UI.
\end{chapterquote}

\newthought{The configuration file} is \code{config/proxy.yml}. Environment variable
interpolation is supported: \code{\$\{VAR\_NAME\}} is replaced with the environment
variable at startup. Secrets (API keys, Redis password) must be provided via
environment variables — never hard-coded in the config file.

\section{Hot Reload Behaviour}
\index{hot reload!configuration}

On \code{SIGHUP}, \ja\ re-reads \code{config/proxy.yml} and applies all changes that
support hot-reload. A \code{config\_reload} event is emitted to the policy audit log
in Redis. The Management UI sends a Redis pub/sub message to trigger reload without
a Unix signal.

\begin{fullwidth}
\begin{table}[h]
\caption{Hot-reload support by configuration section}
\label{tab:hotreload}
\begin{tabular}{@{}llp{6cm}@{}}
\toprule
\textbf{Section} & \textbf{Hot-reloadable} & \textbf{Notes} \\
\midrule
\code{proxy.bind\_host} / \code{bind\_port} & No & Requires process restart \\
\code{proxy.backend\_host} / \code{backend\_port} & Yes & Applies to new connections \\
\code{dial} & Yes & Applies immediately to all new decisions \\
\code{security\_policy} & Yes & New bypass states apply to next connection \\
\code{geoip} & Yes & New country lists apply immediately \\
\code{security} (lists) & Yes & List changes apply to next connection \\
\code{abuseipdb} & Yes & Cache TTL change affects new lookups \\
\code{rdap} & Yes & block\_expansion toggle applies immediately \\
\code{spamhaus} & No (refresh\_interval) & Next trie refresh uses new URL \\
\code{redis.url} / \code{redis.password} & No & Requires process restart \\
\bottomrule
\end{tabular}
\end{table}
\end{fullwidth}

\section{proxy — Core Listener Settings}
\index{configuration!proxy}

\begin{fullwidth}
\begin{table}[h]
\caption{\code{proxy} section — core listener settings}
\label{tab:config-proxy}
\begin{tabular}{@{}lllp{5.5cm}@{}}
\toprule
\textbf{Key} & \textbf{Type} & \textbf{Default} & \textbf{Description} \\
\midrule
\code{bind\_host} & string & \code{"0.0.0.0"} & Listen address for incoming TCP connections from HAProxy \\
\code{bind\_port} & integer & \code{8080} & Listen port. Must match HAProxy backend configuration \\
\code{backend\_host} & string & \textit{required} & Hostname or IP of the protected backend. Must be set \\
\code{backend\_port} & integer & \code{443} & Backend port. Typically 443 for HTTPS \\
\bottomrule
\end{tabular}
\end{table}
\end{fullwidth}

\begin{yamlcode}
proxy:
  bind_host: "0.0.0.0"
  bind_port: 8080
  backend_host: "backend"   # Docker service name or hostname
  backend_port: 443
\end{yamlcode}

\section{dial — Blocking Aggression}
\index{dial}
\index{configuration!dial}

\begin{fullwidth}
\begin{table}[h]
\caption{\code{dial} setting}
\label{tab:config-dial}
\begin{tabular}{@{}lllp{7cm}@{}}
\toprule
\textbf{Key} & \textbf{Type} & \textbf{Default} & \textbf{Description} \\
\midrule
\code{dial} & integer & \code{0} & Blocking aggression: 0 = monitor only (no blocking); 100 = full blocking per configured thresholds. Intermediate values scale blocking proportionally. Raise incrementally. \\
\bottomrule
\end{tabular}
\end{table}
\end{fullwidth}

\begin{warningbox}[Never Deploy with dial=100 on First Install]
Always start at \code{dial: 0}. Observe the score distribution for at least 24 hours
before raising the dial. A premature jump to 100 can block legitimate traffic before
you have validated your threshold configuration.
\end{warningbox}

\section{security\_policy — Bypass Conditions}
\index{configuration!security\_policy}
\index{bypass!configuration}

Each bypass condition has a single \code{enabled} boolean key. All default to \code{true}.

\begin{fullwidth}
\begin{table}[h]
\caption{\code{security\_policy} section — bypass conditions}
\label{tab:config-secpolicy}
\begin{tabular}{@{}llp{3.5cm}p{5.5cm}@{}}
\toprule
\textbf{Key} & \textbf{Default} & \textbf{Effect when true} & \textbf{Risk if disabled} \\
\midrule
\code{alpn\_browser\_bypass.enabled} & \code{true} & h2/h1 ALPN → ALLOW without scoring & Browser traffic scored; false positive risk elevated \\
\code{ja4\_whitelist\_bypass.enabled} & \code{true} & JA4 in whitelist → ALLOW without scoring & Whitelisted fingerprints scored; may be blocked if above threshold \\
\code{mtls\_bypass.enabled} & \code{true} & Valid mTLS cert → ALLOW without scoring & mTLS clients scored; may be blocked if above threshold \\
\code{static\_ip\_allowlist.enabled} & \code{true} & IP in allowlist → ALLOW without scoring & Allowlisted IPs scored \\
\code{ja4\_blacklist\_bypass.enabled} & \code{true} & JA4 in blacklist → TCP RST immediately & Blacklisted fingerprints only blocked if score exceeds dial threshold \\
\code{country\_blacklist\_bypass.enabled} & \code{true} & Country in blacklist → BLOCK immediately & Blacklisted countries only blocked by scorer \\
\code{spamhaus\_bypass.enabled} & \code{true} & Spamhaus DROP match → BLOCK immediately & Spamhaus matches produce RiskSignal(+80) instead of hard block \\
\code{tls\_version\_bypass.enabled} & \code{true} & TLS 1.0/1.1/SSLv3 → TCP RST immediately & Old TLS versions produce risk score contribution instead of hard reject \\
\bottomrule
\end{tabular}
\end{table}
\end{fullwidth}

\begin{yamlcode}
security_policy:
  alpn_browser_bypass:
    enabled: true   # h2/h1 ALPN -> ALLOW without scoring
  ja4_whitelist_bypass:
    enabled: true   # JA4 in whitelist -> ALLOW without scoring
  mtls_bypass:
    enabled: true   # Valid mTLS cert -> ALLOW without scoring
  static_ip_allowlist:
    enabled: true   # IP in static allowlist -> ALLOW without scoring
  ja4_blacklist_bypass:
    enabled: true   # JA4 in blacklist -> immediate RST, no scoring
  country_blacklist_bypass:
    enabled: true   # Country in GeoIP blocklist -> immediate block
  spamhaus_bypass:
    enabled: true   # Spamhaus DROP/EDROP -> immediate block
  tls_version_bypass:
    enabled: true   # TLS 1.0/1.1/SSLv3 -> immediate RST
\end{yamlcode}

\section{geoip — Geographic Controls}
\index{configuration!geoip}
\index{GeoIP}

\begin{fullwidth}
\begin{table}[h]
\caption{\code{geoip} section}
\label{tab:config-geoip}
\begin{tabular}{@{}lllp{5.5cm}@{}}
\toprule
\textbf{Key} & \textbf{Type} & \textbf{Default} & \textbf{Description} \\
\midrule
\code{country\_whitelist\_enabled} & boolean & \code{true} & Enable country allowlist mode. If true, only countries in \code{country\_whitelist} are allowed \\
\code{country\_whitelist} & list & \code{["IE","GB","IM","US"]} & ISO 3166-1 alpha-2 country codes permitted when whitelist mode is active \\
\code{country\_blacklist\_enabled} & boolean & \code{true} & Enable country blacklist mode. Listed countries are immediately blocked \\
\code{country\_blacklist} & list & \code{["KP","RU"]} & ISO 3166-1 alpha-2 country codes to block immediately \\
\bottomrule
\end{tabular}
\end{table}
\end{fullwidth}

\begin{notebox}[Whitelist and Blacklist Together]
Both modes can be active simultaneously. The whitelist check runs first. A country
not on the whitelist (when whitelist mode is active) is blocked independently of
whether it appears on the blacklist. Use whitelist mode for strict geographic restrictions
(allowlist your expected countries), and blacklist mode for blocking specific known-bad
countries when you serve a broad geographic audience.
\end{notebox}

\begin{yamlcode}
geoip:
  country_whitelist_enabled: true
  country_whitelist: ["IE", "GB", "IM", "US"]
  country_blacklist_enabled: true
  country_blacklist: ["KP", "RU"]
\end{yamlcode}

\section{security — Fingerprint Lists and Rate Limits}
\index{configuration!security}

\subsection{Fingerprint Lists}

\begin{fullwidth}
\begin{table}[h]
\caption{\code{security} section — fingerprint lists}
\label{tab:config-fp-lists}
\begin{tabular}{@{}lllp{5.5cm}@{}}
\toprule
\textbf{Key} & \textbf{Type} & \textbf{Default} & \textbf{Description} \\
\midrule
\code{whitelist} & list of strings & \code{[]} & Exact JA4 fingerprints to whitelist. Loaded into Redis \code{ja4:whitelist} and an in-process set \\
\code{whitelist\_patterns} & list of strings & \code{["h2","h1"]} & ALPN patterns that trigger the ALPN browser bypass. Separate from the JA4 whitelist \\
\code{blacklist} & list of strings & \code{[]} & Exact JA4 fingerprints to blacklist. Loaded into Redis \code{ja4:blacklist} \\
\code{fingerprint\_labels} & map & \code{\{\}} & Human-readable names for fingerprints. Used in logs and metrics \\
\code{multi\_strategy\_policy} & string & \code{"majority"} & How to combine rate limit strategy results: \code{any} / \code{majority} / \code{all} \\
\bottomrule
\end{tabular}
\end{table}
\end{fullwidth}

\subsection{Rate Limit Strategies}
\index{rate limiting!configuration}

Three rate-limiting strategies are defined under \code{security.rate\_limit\_strategies}.
Each has an identical structure:

\begin{fullwidth}
\begin{table}[h]
\caption{\code{rate\_limit\_strategies} — per-strategy keys}
\label{tab:config-rl}
\begin{tabular}{@{}lllp{5cm}@{}}
\toprule
\textbf{Key} & \textbf{Type} & \textbf{Default} & \textbf{Description} \\
\midrule
\code{thresholds.suspicious} & integer & — & Connections/second at which the strategy flags as suspicious \\
\code{thresholds.block} & integer & — & Connections/second at which the strategy triggers a block \\
\code{thresholds.ban} & integer & — & Connections/second at which the strategy triggers a permanent ban \\
\code{action} & string & — & Action to take when threshold is crossed: \code{tarpit} or \code{block} \\
\code{ban\_duration} & integer & \code{300} & Initial ban TTL in seconds. Escalates on re-ban \\
\bottomrule
\end{tabular}
\end{table}
\end{fullwidth}

\begin{yamlcode}
security:
  whitelist:
    - "t13d1516h2_8daaf6152771_02713d6af862"  # Chrome 120+
  whitelist_patterns:
    - "h2"
    - "h1"
  blacklist:
    - "t13d190900_9dc949149365_97f8aa674fd9"  # Sliver C2
  fingerprint_labels:
    "t13d1516h2_8daaf6152771_02713d6af862": "Chrome 120+"
  multi_strategy_policy: "majority"
  rate_limit_strategies:
    by_ip:
      thresholds: {suspicious: 50, block: 200, ban: 500}
      action: "tarpit"
      ban_duration: 300
    by_ja4:
      thresholds: {suspicious: 20, block: 100, ban: 200}
      action: "block"
      ban_duration: 300
    by_ip_ja4_pair:
      thresholds: {suspicious: 20, block: 50, ban: 100}
      action: "tarpit"
      ban_duration: 300
\end{yamlcode}

\section{abuseipdb — Reputation Lookup}
\index{configuration!abuseipdb}
\index{AbuseIPDB}

\begin{fullwidth}
\begin{table}[h]
\caption{\code{abuseipdb} section}
\label{tab:config-abuseipdb}
\begin{tabular}{@{}lllp{5.5cm}@{}}
\toprule
\textbf{Key} & \textbf{Type} & \textbf{Default} & \textbf{Description} \\
\midrule
\code{enabled} & boolean & \code{true} & Enable AbuseIPDB lookups \\
\code{api\_key} & string & \textit{required} & AbuseIPDB API key. Must be an env var reference: \code{\$\{ABUSEIPDB\_API\_KEY\}} \\
\code{cache\_ttl} & integer & \code{3600} & Seconds to cache AbuseIPDB responses in Redis \\
\code{min\_score\_to\_signal} & integer & \code{50} & Minimum AbuseIPDB score (0–100) that produces a RiskSignal. Never hard-block below 50 \\
\bottomrule
\end{tabular}
\end{table}
\end{fullwidth}

\begin{warningbox}[AbuseIPDB Score 50 Floor]
The \code{min\_score\_to\_signal} floor exists because AbuseIPDB scores below 50 are
common for shared infrastructure (NAT gateways, VPNs, corporate egress). Hard-blocking
based on scores below 50 has an unacceptably high false positive rate. The default of
50 should not be lowered without evaluating the false positive impact on your traffic.
\end{warningbox}

\begin{yamlcode}
abuseipdb:
  enabled: true
  api_key: "${ABUSEIPDB_API_KEY}"
  cache_ttl: 3600
  min_score_to_signal: 50
\end{yamlcode}

\section{rdap — WHOIS/RDAP Org Reputation}
\index{configuration!rdap}
\index{RDAP}

\begin{fullwidth}
\begin{table}[h]
\caption{\code{rdap} section}
\label{tab:config-rdap}
\begin{tabular}{@{}lllp{5.5cm}@{}}
\toprule
\textbf{Key} & \textbf{Type} & \textbf{Default} & \textbf{Description} \\
\midrule
\code{enabled} & boolean & \code{true} & Enable RDAP/WHOIS enrichment lookups \\
\code{block\_expansion} & boolean & \code{false} & Automatically block the entire /24 (IPv4) or /48 (IPv6) of a blocked IP. Must be explicitly enabled \\
\code{max\_expansion} & integer & \code{24} & Maximum prefix length for block expansion. Never set below 24 for IPv4 \\
\bottomrule
\end{tabular}
\end{table}
\end{fullwidth}

\begin{dangerbox}[Block Expansion Off by Default]
\code{block\_expansion: true} will cause \ja\ to block the entire /24 subnet when a
single IP is blocked. On residential ISPs and shared hosting, an entire /24 may serve
thousands of unrelated users. Enable only after confirming that the blocked IPs are
from dedicated attacker infrastructure with no residential or business customers.
\end{dangerbox}

\begin{yamlcode}
rdap:
  enabled: true
  block_expansion: false      # secops must explicitly opt in
  max_expansion: 24           # never beyond /24 IPv4, /48 IPv6
\end{yamlcode}

\section{spamhaus — DROP/EDROP Blocklists}
\index{configuration!spamhaus}
\index{Spamhaus}

\begin{fullwidth}
\begin{table}[h]
\caption{\code{spamhaus} section}
\label{tab:config-spamhaus}
\begin{tabular}{@{}lllp{5.5cm}@{}}
\toprule
\textbf{Key} & \textbf{Type} & \textbf{Default} & \textbf{Description} \\
\midrule
\code{drop\_list\_url} & string & \footnotesize\code{https://www.spamhaus.org/drop/drop.txt} & URL to the Spamhaus DROP list \\
\code{edrop\_list\_url} & string & \footnotesize\code{https://www.spamhaus.org/drop/edrop.txt} & URL to the Spamhaus EDROP list \\
\code{refresh\_interval} & integer & \code{3600} & Seconds between trie refresh cycles \\
\bottomrule
\end{tabular}
\end{table}
\end{fullwidth}

\begin{yamlcode}
spamhaus:
  drop_list_url: "https://www.spamhaus.org/drop/drop.txt"
  edrop_list_url: "https://www.spamhaus.org/drop/edrop.txt"
  refresh_interval: 3600
\end{yamlcode}

\section{redis — Connection Settings}
\index{configuration!redis}
\index{Redis!connection}

\begin{fullwidth}
\begin{table}[h]
\caption{\code{redis} section}
\label{tab:config-redis}
\begin{tabular}{@{}lllp{5.5cm}@{}}
\toprule
\textbf{Key} & \textbf{Type} & \textbf{Default} & \textbf{Description} \\
\midrule
\code{url} & string & \code{"redis://redis:6379"} & Redis connection URL. For Unix socket: \code{unix:///var/run/redis/redis.sock} \\
\code{password} & string & \textit{optional} & Redis AUTH password. Use env var: \code{\$\{REDIS\_PASSWORD\}} \\
\bottomrule
\end{tabular}
\end{table}
\end{fullwidth}

\begin{notebox}[Unix Domain Socket Performance]
For deployments where \ja\ and Redis are on the same host, using a Unix domain socket
for the Redis URL reduces latency by approximately 30–40\% compared to TCP loopback.
Set \code{url: "unix:///var/run/redis/redis.sock"} and ensure the socket path is
accessible inside the container.
\end{notebox}

\begin{yamlcode}
redis:
  url: "redis://redis:6379"
  password: "${REDIS_PASSWORD}"
\end{yamlcode}

\section{Complete Configuration Example}
\index{configuration!example}

\begin{yamlcode}
# config/proxy.yml — complete reference configuration

proxy:
  bind_host: "0.0.0.0"
  bind_port: 8080
  backend_host: "backend"
  backend_port: 443

# 0 = monitor only. Raise incrementally after validating score distribution.
dial: 0

security_policy:
  alpn_browser_bypass:
    enabled: true
  ja4_whitelist_bypass:
    enabled: true
  mtls_bypass:
    enabled: true
  static_ip_allowlist:
    enabled: true
  ja4_blacklist_bypass:
    enabled: true
  country_blacklist_bypass:
    enabled: true
  spamhaus_bypass:
    enabled: true
  tls_version_bypass:
    enabled: true

geoip:
  country_whitelist_enabled: true
  country_whitelist: ["IE", "GB", "IM", "US"]
  country_blacklist_enabled: true
  country_blacklist: ["KP", "RU"]

security:
  whitelist:
    - "t13d1516h2_8daaf6152771_02713d6af862"  # Chrome 120+
  whitelist_patterns:
    - "h2"
    - "h1"
  blacklist:
    - "t13d190900_9dc949149365_97f8aa674fd9"  # Sliver C2
  fingerprint_labels:
    "t13d1516h2_8daaf6152771_02713d6af862": "Chrome 120+"
  multi_strategy_policy: "majority"
  rate_limit_strategies:
    by_ip:
      thresholds: {suspicious: 50, block: 200, ban: 500}
      action: "tarpit"
      ban_duration: 300
    by_ja4:
      thresholds: {suspicious: 20, block: 100, ban: 200}
      action: "block"
      ban_duration: 300
    by_ip_ja4_pair:
      thresholds: {suspicious: 20, block: 50, ban: 100}
      action: "tarpit"
      ban_duration: 300

abuseipdb:
  enabled: true
  api_key: "${ABUSEIPDB_API_KEY}"
  cache_ttl: 3600
  min_score_to_signal: 50

rdap:
  enabled: true
  block_expansion: false
  max_expansion: 24

spamhaus:
  drop_list_url: "https://www.spamhaus.org/drop/drop.txt"
  edrop_list_url: "https://www.spamhaus.org/drop/edrop.txt"
  refresh_interval: 3600

redis:
  url: "redis://redis:6379"
  password: "${REDIS_PASSWORD}"
\end{yamlcode}
